
本书是一套两卷本专著。本系列自足地阐述开弦与闭弦两种情形下的基础\termFloerHomology{}，并系统介绍它在哈密顿动力学和辛拓扑问题中的应用。这两卷书的基本研究对象，是\termLagrangianSubmanifold{}的几何、\termHamiltonianDiffeomorphism{}的动力学，以及二者在辛拓扑中的相互作用。

辛几何中的经典达布定理揭示了辛变换群的\emph{柔性}。另一方面，米哈伊尔·格罗莫夫（Mikhail Gromov）与雅科夫·埃利亚什伯格（Yakov Eliashberg）的著名定理 (El87) 揭示了辛变换微妙的\emph{刚性}：\emph{由辛微分同胚构成的子群 $\Symp(M,\omega)$ 在 $\Diff(M)$ 中关于 $C^0$ 拓扑是闭的}。这表明辛拓扑的研究既微妙又有趣。埃利亚什伯格的定理依赖于某个版本的非挤压定理，例如格罗莫夫 (Gr85) 运用\termPseudoholomorphicCurve{}工具所证明的版本。除了埃利亚什伯格对这一非挤压结果最初给出的组合证明之外，伊瓦尔·埃克兰（Ivar Ekeland）与赫尔穆特·霍费尔（Helmut Hofer）(EkH89) 还利用哈密顿系统的经典变分方法给出了另一个证明。辛几何的这两个侧面——\termPseudoholomorphicCurve{}的分析与哈密顿动力学——之间的相互作用，自安德烈亚斯·弗洛尔（Andreas Floer）对其半无限维同调理论（如今称为\emph{\termFloerHomology{}}理论）的开创性工作以来，一直是辛拓扑发展的主要推动力。

哈密顿方程 $\dot{x}=X_H(t,x)$ 源自哈密顿力学；自约瑟夫-路易·拉格朗日（Joseph-Louis Lagrange）、威廉·罗恩·哈密顿（William Rowan Hamilton）、卡尔·古斯塔夫·雅各布·雅可比（Carl Gustav Jacob Jacobi）、昂利·庞加莱（Henri Poincaré）等人的时代起，对其动力学的研究便一直是物理学中的基本主题。为了理解它的动力学并求出方程的显式解，人们发展了许多数学工具。研究这些问题的一个关键工具是\emph{最小作用量原理}：哈密顿方程的一个解对应于某个作用量泛函的临界点。在这一变分原理中，对一般\termSymplecticManifold{}上的方程 $\dot{x}=X_H(t,x)$，最重要的两类边界条件如下：一类是\emph{周期}边界条件 $\gamma(0)=\gamma(1)$；另一类是\emph{拉格朗日}边界条件
\[
  \gamma(0)\in L_0,\qquad \gamma(1)\in L_1,
\]
其中 $(L_0,L_1)$ 是给定的一对\termLagrangianSubmanifold{}。若子流形 $i:L\hookrightarrow(M,\omega)$ 满足 $i^*\omega=0$ 且 $\dim L=\frac12\dim M$，则称它为拉格朗日的。这取代了经典力学中的\emph{两点}边界条件。

若\termSymplecticManifold{} $(M,\omega)$ 的微分同胚 $\phi$ 是某个（依赖时间的）哈密顿量 $H$ 所对应方程 $\dot{x}=X_H(t,x)$ 的时间一映射，则称 $\phi$ 为\termHamiltonianDiffeomorphism{}。所有这类微分同胚组成的集合记为 $\Ham(M,\omega)$；它构成 $\Symp(M,\omega)$ 的一个子群。然而，在作者看来，之所以研究\termHamiltonianDiffeomorphism{}，纯粹是历史上的偶然，因为 $\Ham(M,\omega)$ 的定义并非先验自然的。例如，与 $\Symp(M,\omega)$ 是辛结构 $\omega$ 的自同构群不同，$\Ham(M,\omega)$ 并不是光滑流形 $M$ 所带任何几何结构的结构群（至少迄今尚未发现这样的结构）。数学家对 $\Ham(M,\omega)$ 的兴趣，很大程度上由著名的阿诺尔德猜想 (Ar65) 所推动；弗洛尔正是在尝试证明这一猜想的过程中创立了\termFloerHomology{} (Fl88b, Fl89b)。

自二十世纪八十年代末\termFloerHomology{}出现以来，它一直在辛拓扑的发展中扮演基础性角色。（低维拓扑中也有一个平行的概念，本书两卷不予涉及。对此感兴趣的读者，我们推荐弗洛尔的原始论文 (Fl88c) 以及彼得·克朗海默（Peter Kronheimer）与托马什·姆罗夫卡（Tomasz Mrowka）的杰作 (KM07)。）由于严格定义中牵涉许多技术细节，尤其是在拉格朗日交的\termFloerHomology{}（或“开弦”情形）中，这一主题对初学的研究生和来自其他数学领域的研究人员而言颇难入门。部分原因在于，现有文献尚未系统解释辛拓扑中的问题、分析细节，以及把整个弗洛尔理论所包含的工具应用起来所涉及的技术。与此同时，深谷贤治（Kenji Fukaya）对\termFloerHomology{}的范畴化，即他把一个 $A_\infty$ 范畴引入辛几何（现称深谷范畴），以及马克西姆·孔采维奇（Maxim Kontsevich）的同调镜像对称纲领，继而还有物理学中 D-膜开弦理论的发展，都极大地推进了弗洛尔理论，并且不仅吸引了传统辛几何学家和拓扑学家，也吸引了其他数学家和物理学家的广泛关注。此外，把辛几何思想应用于二维（保面积）动力系统中各类问题的研究也已相当丰富。

我们撰写这两卷书，是希望在一定程度上缓解当前的这些困难。为实现这一目标，我们更集中于作者较为熟悉的弗洛尔理论基础材料及其对辛拓扑中各种问题的应用；并在假定读者已经掌握（一阶）偏微分方程的基础椭圆理论、研究生二年级程度的微分几何和研究生一年级程度的代数拓扑的前提下，力求给出完整的分析细节。此外，我们还尝试从经典拉格朗日--哈密顿变分原理与哈密顿力学的历史背景出发，为弗洛尔理论中出现的各种构造提供动机。本书选题多少带有一定偏向，一方面受限于作者的知识范围和把握程度，另一方面也源于作者试图避免与现有辛拓扑文献过多重叠。在诸多著作中，我们特别列出以下三部专著，并将本书两卷与它们作比较：

\begin{enumerate}[label=(\arabic*)]
  \item \emph{J-Holomorphic Curves and Symplectic Topology}, 杜萨·麦克达夫（Dusa McDuff）、迪特马尔·萨拉蒙（Dietmar Salamon）, 2004.
  \item \emph{Fukaya Categories and Picard--Lefschetz Theory}, 保罗·塞德尔（Paul Seidel）, 2008.
  \item \emph{Lagrangian Intersection Floer Theory: Anomaly and Obstruction}, volumes I \& II, 深谷贤治、吴龙根、太田啓史（Hiroshi Ohta）、小野薫（Kaoru Ono）, 2009.
\end{enumerate}

（另有米歇尔·奥丹（Michèle Audin）与米哈伊·达米安（Mihai Damian）较新的一部专著 (AD14)，原以法文写成，后来译为英文。）

首先，这两卷书的第 2、3 部分可以看作研究生或博士后阅读第 (3) 本书 (FOOO09) 的预备材料，因为第二卷中拉格朗日弗洛尔理论的离壳设置预设了非平凡瞬子，即非常值全纯圆盘或球面的存在。不过，我们基本限于单调情形，并避开成熟完备的\termFloerHomology{}障碍--形变理论，因为后者不可避免地会牵涉 $A_\infty$ 结构理论以及仓西结构等虚拟模空间技术的抽象扰动理论，而这些超出了本书两卷的范围。所幸第 (2) 本书 (Se08) 和第 (3) 本书 (FOOO09) 涵盖了理论的这一重要方面，因此我们强烈建议读者参考它们。除余切丛情形以外，我们也基本不对正合\termLagrangianSubmanifold{}的弗洛尔理论作广泛讨论，因为塞德尔的著作 (Se08) 已在正合辛几何的背景下对弗洛尔理论和深谷范畴作了深入研究，我们无可增益之处。第 2 部分关于基础\termPseudoholomorphicCurve{}理论的材料，与麦克达夫和萨拉蒙的第 (1) 本书 (MSa04) 第 1--6 章有许多重叠。然而，我们阐述这些材料的方式与 (MSa04) 颇为不同。

例如，我们从一开始便处理任意亏格且带边界的\termPseudoholomorphicCurve{}，并统一处理闭情形与开情形，例如弱解的正则性理论和可去奇点定理就是如此。我们还把横截性问题放在紧性问题之后讨论；当读者希望超越半正情形而采用仓西结构和抽象扰动理论的技术时，这样的顺序似乎更为合适。在阐述\termPseudoholomorphicCurve{}理论时，我们还特别关注另外两点。其一，是使相关几何分析更接近黎曼几何中更为标准的几何分析风格：只要可能，就强调利用与殆凯勒性质相伴的典范联络来作张量计算。这样，我们通过纯粹的张量计算和魏岑伯克公式的应用，推导出相关的 $W^{k,p}$ 强制估计，特别是具有诺伊曼边界条件的最优 $W^{2,2}$ 估计。其二，是使\termPseudoholomorphicCurve{}的形变理论接近（可积）凯勒流形上全纯曲线的形变理论。我们希望这样的阐述方式能把读者群从传统的辛几何学家拓展到其他数学领域的研究生和研究人员，使他们能够更容易接触辛拓扑及相关领域的重要进展。

下面简要概述这两卷书各部分的内容。第一卷由第 1、2 部分组成。

第 1 部分从经典力学中的拉格朗日与哈密顿经典变分原理出发，介绍辛几何，并引入辛几何中的主要概念，即\termLagrangianSubmanifold{}、\termHamiltonianDiffeomorphism{}和辛纤维化。它还介绍\termHamiltonianDiffeomorphism{}的霍费尔几何。本部分最后证明格罗莫夫--埃利亚什伯格 $C^0$-辛刚性定理 (El87)，并介绍连续哈密顿动力学以及斯特凡·穆勒（Stefan Müller）与吴龙根所提出的哈密顿同胚概念 (OhM07)。

第 2 部分对\termPseudoholomorphicCurve{}及其模空间的分析作了大体自足的阐述。我们试图在\emph{离壳设置}下给出非线性柯西--黎曼算子 $\overline\partial_J$ 的先验椭圆估计的最优形式。为此，我们强调使用\termAlmostKahlerManifold{} $(M,\omega,J)$ 的典范联络。我们在分析处理上的另一项新意，是对满足全实边界条件的 $J$-全纯曲线弱解（叶如钢（Rugang Ye）(Ye94) 意义下）的边界正则性定理给出完整证明。据我们所知，现有文献此前尚未给出这个正则性证明。我们还沿用深谷贤治与小野薫 (FOn99) 的方法，完整证明稳定映射模空间的紧性。本部分最后说明：如何将（受扰动的）\termPseudoholomorphicCurve{}模空间的紧性--配边分析与少量辛拓扑数据相结合，从而证明辛拓扑中的两个基本定理，即格罗莫夫非挤压定理以及（对驯服\termSymplecticManifold{}而言）$\Ham(M,\omega)$ 上\termHoferNorm{}的非退化性。

第二卷由第 3、4 部分组成。第 3 部分介绍限于单调\termLagrangianSubmanifold{}特殊情形的拉格朗日\termFloerHomology{}。它先概述余切丛上的拉格朗日交\termFloerHomology{}，并引入近期拉格朗日交弗洛尔理论中的所有主要研究对象，而不过多深入技术细节。然后，它解释弗洛尔模空间的紧化；相关细节在文献中往往含混不清。本部分最后构造一个谱序列，研究可位移\termLagrangianSubmanifold{}的马斯洛夫类障碍，并介绍波尔捷罗维奇关于 $\Ham(S^2)$ 的霍费尔直径的定理。

第 4 部分介绍哈密顿\termFloerHomology{}，解释谱不变量的完整构造及其各种应用。这些应用包括谱范数的构造、迈克尔·厄舍（Michael Usher）对霍费尔几何中极小性猜想的证明，以及最优能量--容量不等式。特别地，本部分完整、自足地阐述恩托夫--波尔捷罗维奇的谱拟同态构造及其相伴的辛拟态。本部分最后进一步讨论拓扑哈密顿流，以及它们与二维保面积同胚几何之间的关系。

阅读这两卷书所需的预备知识随各部分而异。标准的研究生一年级可微流形课程，再加上一点纤维丛理论知识，应足以阅读第 1 部分。然而，第 2 部分是全书技术性最强的部分，处理\termPseudoholomorphicCurve{}的一般理论、其模空间及其稳定映射紧化；它假定读者掌握黎曼几何的基本语言（例如迈克尔·斯皮瓦克（Michael Spivak）(Spi79) 第 1、2 卷所述内容）、基础泛函分析（例如索伯列夫嵌入、雷利希紧性等）、椭圆型（一阶）偏微分方程以及研究生一年级代数拓扑。第 3、4 部分处理开弦和闭弦情形下\termFloerHomology{}的主要主题，依赖第 1、2 部分的材料，但本身应当能够独立阅读。已经熟悉基础辛几何和\termPseudoholomorphicCurve{}分析的读者，可以直接阅读第 3、4 部分。本书可以用作研究生教材，介绍格罗莫夫与弗洛尔研究现代辛拓扑的分析方法。希望深入了解辛拓扑和镜像对称各方面内容的读者，强烈建议依兴趣另行阅读上列第 (1)--(3) 本书。

在结束前言之际，作者愿回顾自己的个人经历和看法；这些未必为他人所共有，但作者希望它们能帮助读者理解这两卷书为何形成目前的面貌。辛拓扑这一概念肇始于埃利亚什伯格与格罗莫夫著名的辛刚性定理。埃利亚什伯格最初的证明基于辛微分同胚的某种 $C^0$ 型不变量的存在性，该不变量度量辛向量空间中区域的大小。这类不变量的存在性最初由格罗莫夫建立，它是格罗莫夫借助\termPseudoholomorphicCurve{}分析方法证明的基本\emph{非挤压定理}的一个推论。随着格罗莫夫所发展的\termPseudoholomorphicCurve{}方法问世，以及弗洛尔随后创立导出\termFloerHomology{}的椭圆型莫尔斯理论，辛几何的面貌发生了巨变。许多过去无法处理的辛几何问题借助\termPseudoholomorphicCurve{}技术得到解决，辛拓扑这一概念也逐渐成形。

作者以目前方式安排本书结构，主要受两个因素影响。第一个因素关系到第一卷如何处理分析材料。\termPseudoholomorphicCurve{}方法刚出现时所伴随的困难（或者说令人兴奋之处），在于它涉及一种数学学科，其性质与当时主流辛几何学家所使用的数学类型大为不同。因此，作者感到，格罗莫夫的论文问世前后的辛几何之间产生了某种断裂；其中呈现的分析，也与黎曼几何中的几何分析学家通常处理这类问题的方式颇不相同。例如，张量计算并未像在黎曼几何中那样受到强调。此外，从作者个人经历来看，无论作为研究生，还是作为刚开始从事辛拓扑研究的研究人员，进入格罗莫夫--威腾--弗洛尔理论都有两道障碍。第一道，是必须克服对稳定曲线的德利涅--芒福德模空间 $\mathcal M_{g,n}$ 之类抽象代数几何材料的某种畏惧；第二道，则是必须从相关数学文献的原始资料中吸收研究\termPseudoholomorphicCurve{}映射模空间所需的大量分析材料，而那些资料的细节往往过于简略。事实证明，从几何分析学家的观点看，许多这样的细节在某种意义上都是标准的，并且可以利用黎曼几何的标准张量方法更有效地处理。

第二个因素关系到本书如何呈现弗洛尔理论。依作者的个人经验，同时在闭弦与开弦情形下学习弗洛尔理论似乎最为有效。哈密顿动力学中的问题往往要通过拉格朗日交几何中的相应问题来解决。因此，作者同时呈现闭弦与开弦情形下的弗洛尔理论。尽管\termPseudoholomorphicCurve{}的技术性分析细节在闭弦与开弦情形下本质相同，但对黎曼曲面的闭情形与开情形，相关\termPseudoholomorphicCurve{}模空间的几何却不相同。这一区别使拉格朗日交弗洛尔理论与哈密顿不动点弗洛尔理论大为不同。

\section*{致谢}

本书受惠于许多人，他们给予的宝贵帮助无论怎样强调都不为过。这两卷书的写作计划始于 2002--2004 年间作者在威斯康星大学麦迪逊分校开设的一门持续三个学期的辛几何与\termPseudoholomorphicCurve{}课程，以及 2004--2005 年作者在斯坦福大学休假访问期间开设的另一门持续一个学季的拉格朗日\termFloerHomology{}课程。

作者在加州大学伯克利分校师从艾伦·温斯坦（Alan Weinstein）攻读研究生期间，从他那里学到了基础辛几何的大部分材料。特别是，第一卷第 1 部分的大部分内容都以作者为温斯坦于 1987--1988 年开设的全年辛几何课程所记的课堂笔记为基础。在整个第 1 部分中，人们很容易看到温斯坦的数学所产生的广泛影响。作者愿借此机会，衷心感谢他在自己整个研究生学习期间给予的宝贵支持和鼓励。

作者还感谢埃利亚什伯格（原文称 Yasha Eliashberg），使作者得以在斯坦福大学休假访问，并给予作者开设弗洛尔同调课程的机会。作者也十分感谢彼得·施佩特（Peter Spaeth）与 Cheol-Hyun Cho；他们为作者 2002--2004 年课程所作的课堂笔记成为第一卷的基础。作者还感谢王兵（Bing Wang）出色地录入第一卷初稿；若非这项帮助，这两卷书能否问世将大成疑问。还要感谢 Dongning Wang 与鲍尔考（Erkao Bao）分别校阅第 2、4 部分的早期版本，并感谢王蕊（Rui Wang）与 Yoosik Kim 校阅第 1、2 部分的早期版本。

作者也感谢多年的朋友和合作者深谷贤治、太田啓史与小野薫，感谢他们在拉格朗日交弗洛尔理论及其应用方面的合作。第 3 部分的许多材料都以我们的合作为基础。特别感谢深谷贤治在整个合作过程中展现的杰出领导力和数学远见，作者从中学到了许多数学。

作者感谢迈克尔·恩托夫（Michael Entov）与列昂尼德·波尔捷罗维奇（Leonid Polterovich）对辛拟态和拟同态作出的惊人构造；作者相信，这一构造开始揭示\termFloerHomology{}与谱不变量之分析构造背后真正的辛本性。作者也感谢克洛德·维泰博（Claude Viterbo）、迈克尔·厄舍（Michael Usher）、列夫·布霍夫斯基（Lev Buhovsky）与索布汉·赛法迪尼（Sobhan Seyfaddini；原书作 Soban Seyfaddini）的重要数学工作；这些工作与作者较近期关于谱不变量和拓扑哈密顿动力学的研究直接相关，作者对此深为感激。

作者还感谢在本书写作期间曾停留过的各个机构。特别感谢威斯康星大学麦迪逊分校；作者任教于该校期间完成了这两卷书的大部分写作。作者也感谢韩国高等研究院 (KIAS)、国家数学科学研究所 (NIMS) 与首尔大学，它们在作者停留期间提供了经费支持和优良的研究环境。

最后但同样重要的是，作者毫无保留地感谢已故的弗洛尔；他惊人的数学工作主导了作者的整个数学工作。

在本书准备期间，作者部分得到美国 NSF 的一项资助，并得到韩国 IBS 项目 \# IBS-R003-D1 的支持。

\section*{约定一览}

对于哈密顿向量场、作用量泛函以及哈密顿弗洛尔理论和谱不变量、恩托夫--波尔捷罗维奇卡拉比拟同态构造中出现的其他对象，我们遵循 (Oh05c, Oh09a, Oh10) 的约定。它们在若干方面不同于例如 (EnP03, EnP06, EnP09) 所采用的约定。

\begin{enumerate}[label=(\arabic*)]
  \item 余切丛 $T^*N$ 上的典范\termSymplecticForm{} $\omega_0$ 定义为
  \[
    \omega_0=-d\Theta=\sum_{i=1}^n dq^i\wedge dp_i,
  \]
  其中 $\Theta$ 是刘维尔 1-形式；在由 $N$ 上坐标系
  $(q^1,\ldots,q^n)$ 诱导的典范坐标
  $(q^1,\ldots,q^n,\allowbreak p_1,\ldots,p_n)$ 中，
  它由 $\Theta=\sum_{i=1}^n p_i\,dq^i$ 给出。

  \item \termSymplecticManifold{} $(M,\omega)$ 上的哈密顿向量场 $X_H$ 由 $dH=\omega(X_H,\cdot)$ 定义。

  \item 作用量泛函 $\mathcal A_H:\widetilde{\mathcal L}_0(M)\to\mathbb R$ 定义为
  \[
    \mathcal A_H([\gamma,w])=-\int w^*\omega-\int_0^1 H(t,\gamma(t))\,dt.
  \]

  \item 特别地，在 $T^*N$ 的道路空间上，$\mathcal A_H$ 化为经典哈密顿作用量泛函
  \[
    \int_\gamma p\,dq-\int_0^1 H(t,\gamma(t))\,dt,
  \]
  这与经典力学文献中的标准定义\emph{一致}。

  \item 若\termAlmostComplexStructure{}满足 $\omega(X,JX)\geq 0$，则称它为 $J$-正的；若双线性形式 $\omega(\cdot,J\cdot)$ 定义了一个黎曼度量，则称它为 $J$-相容的。

  \item 注意 $\mathbb R^{2n}\cong\mathbb C^n$ 带有三个典范双线性形式：\termSymplecticForm{} $\omega_0$、\emph{欧氏内积} $g$ 和埃尔米特内积 $\langle\cdot,\cdot\rangle$。我们规定埃尔米特内积对第一个变元复线性、对第二个变元反复线性。三者之间关系的约定为
  \[
    \langle\cdot,\cdot\rangle=g(\cdot,\cdot)-i\omega_0(\cdot,\cdot).
  \]
\end{enumerate}

\section*{与恩托夫--波尔捷罗维奇约定 (EnP09) 的比较}

在 (EnP03, EnP06, EnP09) 中，恩托夫与波尔捷罗维奇所用的符号约定不同于 (Oh05c) 和本书。将我们的约定与 (EnP09) 的约定相比较，差异如下。

\begin{enumerate}[label=(\arabic*)]
  \item 在他们的约定中，$T^*N$ 上的典范\termSymplecticForm{}是
  \[
    \widetilde\omega_0\coloneqq d\Theta=\sum_{i=1}^n dp_i\wedge dq^i.
  \]

  \item (EnP09) 中的哈密顿向量场定义为 $dH=\omega(\cdot,X_H)$。因而，把 $H$ 换成 $-H$ 后，对给定\termSymplecticManifold{} $(M,\omega)$，在两种约定下，相应哈密顿向量场的周期解具有相同的闭环集合。

  \item (1) 与 (2) 相结合，使得余切丛上函数 $H=H(t,q,p)$ 所对应的哈密顿向量场在两种约定中给出\emph{同一个}向量场。特别地，相空间 $\mathbb R^{2n}$ 上关于典范坐标 $(q^1,\ldots,q^n,p_1,\ldots,p_n)$ 的经典哈密顿向量场为
  \[
    X_H=\sum_{i=1}^n\left(
      \frac{\partial H}{\partial p_i}\frac{\partial}{\partial q^i}
      -\frac{\partial H}{\partial q^i}\frac{\partial}{\partial p_i}
    \right).
  \]

  \item (EnP03) 与 (EnP09) 把作用量泛函定义为
  \begin{equation}\tag{0.0.1}\label{eq:enp-action}
    -\int w^*\omega+\int_0^1 H(t,\gamma(t))\,dt.
  \end{equation}
  为在下文比较两种定义，我们把式~\eqref{eq:enp-action} 所给的定义记为 $\widetilde{\mathcal A}_H([\gamma,w])$。

  \item 特别地，在 $T^*N$ 的道路空间上，$\widetilde{\mathcal A}_H$ 化为
  \[
    -\int_\gamma p\,dq+\int_0^1 H(t,\gamma(t))\,dt,
  \]
  它是经典力学文献（例如 (Ar89) 和 (Go80)）中哈密顿作用量泛函标准定义的\emph{负值}。

  \item 由于这两种约定采用同一个相伴殆凯勒度量 $\omega(\cdot,J\cdot)$，相应的受扰动柯西--黎曼方程具有完全相同的形式；特别地，哈密顿向量场前面的符号相同。

  \item 此外，恩托夫与波尔捷罗维奇 (EnP03, EnP06) 用 $c(a,H)$ 表示谱数，其中 $a$ 是量子\emph{同调}类；而我们用 $a$ 表示量子上同调类。二者的比较如下：
  \begin{equation}\tag{0.0.2}\label{eq:enp-spectral-comparison}
    \rho(H;a)=c(a^\flat;\widetilde H)=c(a^\flat;\overline H),
  \end{equation}
  其中 $a^\flat$ 是与上同调类 $a$ 对偶的同调类，而 $\overline H$ 是 $H$ 的逆哈密顿量，由
  \begin{equation}\tag{0.0.3}\label{eq:inverse-hamiltonian}
    \overline H(t,x)=-H(t,\phi_H^t(x))
  \end{equation}
  给出。

  \item 在他们的约定中，$\mathbb R^{2n}$ 上三个双线性形式的关系为
  \[
    \langle\cdot,\cdot\rangle=g(\cdot,\cdot)+i\widetilde\omega_0(\cdot,\cdot),
  \]
  此时该内积对第二个变元复线性、对第一个变元反复线性。
\end{enumerate}

理解了这些之后，就可以把 (EnP03, EnP06) 中的每个陈述翻译成用我们的记号表达的陈述。
